% GENERATED FILE. Edit canonical lessons, then run npm run generate:books.
% canonical-source-hash: fnv1a64:5f9f5f42825bbf6f
% canonical-lessons: ZH-C48-dian, ZH-W48-dian, ZH-C48-fen, ZH-C48-ban, ZH-C48-ji, ZH-C48-deng

\chapter{Telling the Time}
\label{ch:zh-telling-the-time}
\hlchaptermodality{48}

\begin{chapteropening}
\textbf{By the end of this chapter:} \emph{I can ask the time and say it to the half hour and the minute.}

Complete the last lesson of chapter 48: I can ask the time and say it to the half hour and the minute.
\end{chapteropening}

\section[\texorpdfstring{\zh{点} (diǎn)}{点 (diǎn)}]{\zh{点} (diǎn) — o'clock; dot}

\label{lesson:ZH-C48-dian}



\begin{quote}
Before the new one: say the Mandarin for this year, then the Mandarin for date.
\end{quote}

\subsection*{You'll want to know: \zh{点}}
\textbf{\zh{点}} — \emph{diǎn} — \textquotedblleft{}o'clock; dot\textquotedblright{}.

\begin{quote}
\textbf{\zh{我们明天八点去}} \emph{wǒmen míngtiān bā diǎn qù} — We're going at eight tomorrow.
\end{quote}

The new character, in pieces: \zh{点} — a small vertical and bar on top of \zh{口}, standing over four dots. The next lesson puts it in your hand.

\begin{cousinweb}
\textbf{\zh{点}} (\emph{diǎn}) is inherited from Middle Chinese \emph{temX} (Guangyun, in Baxter's transcription), glossed 'a dot made with the brush'. Its use for the hours is not given there, so how it took that job is not settled.
\end{cousinweb}

\begin{grammarlens}[title={What you've built}]
One more word for times and partings.
\end{grammarlens}

\subsection*{Your turn}
\begin{itemize}
\raggedright
  \item \emph{Say it:} \emph{diǎn}
  \item \emph{Say it:} \emph{diǎn}, once more
  \item \emph{Say it:} \emph{wǒmen míngtiān bā diǎn qù} — We're going at eight tomorrow.
  \item \emph{Recall:} say the Mandarin for this year, then the Mandarin for date, then say \emph{diǎn} again
\end{itemize}

\begin{tcolorbox}[breakable,colback=teal!4,colframe=teal!35!black,title={Before you move on}]
What does \zh{点} mean? (\textquotedblleft{}O'clock; dot\textquotedblright{}.) Say it once more.
\end{tcolorbox}
\section[\texorpdfstring{\zh{点} (diǎn)}{点 (diǎn)}]{\zh{点} — a mark on \zh{口}, over four dots}

\label{lesson:ZH-W48-dian}



\begin{quote}
The word first, then the character inside it.

\begin{itemize}
\raggedright
  \item \emph{Recall:} say \emph{diǎn} — \textbf{R1}, one lesson back, before it decays
\end{itemize}

\begin{quote}
\zh{口} \zh{明}
\end{quote}

\emph{Read it:} each one aloud, with nothing said first — \emph{kǒu}, \emph{míng}
\end{quote}

\subsection*{Script — a mark on \zh{口}, over four dots}
\hlblockfigure{\detokenize{figures/ZH-W48-dian-filmstrip.pdf}}{How \zh{点} is written, stroke by stroke}

\begin{quote}
\zh{点}
\end{quote}

\textbf{Three pieces, top to bottom.}

\textbf{Two strokes on top.} A short vertical, then a short horizontal from its middle out to the right.

\textbf{\zh{口}, three strokes}, wide, under them.

\textbf{Four dots} along the bottom, left to right; the first leans left, the other three lean right.

2 + 3 + 4 = nine.

\subsection*{Your turn}
\begin{itemize}
\raggedright
  \item \emph{Write it:} \zh{点}, all nine, with the model in front of you
  \item \emph{Cover:} write it again from memory, and count the pieces, not the strokes
\end{itemize}

\begin{tcolorbox}[breakable,colback=teal!4,colframe=teal!35!black,title={Before you move on}]
How many pieces in \zh{点}? (\textbf{Three}.) Which character you know is in the middle? (\textbf{\zh{口}}.) How many dots at the bottom? (\textbf{Four}.)

Source: \href{https://raw.githubusercontent.com/chanind/hanzi-writer-data/68d10a4b21150cae5e1ebbd223eed289cf32d90c/data/\%E7\%82\%B9.json}{Hanzi Writer Data \zh{点}, pinned stroke paths and medians}.
\end{tcolorbox}
\section[fēn]{fēn — minute; to divide}

\label{lesson:ZH-C48-fen}



\begin{quote}
Before the new one: say the Mandarin for date, then the Mandarin for o'clock.
\end{quote}

\subsection*{You'll want to know: fēn}
\emph{fēn} — \textquotedblleft{}minute; to divide\textquotedblright{}.

\begin{quote}
\emph{bā diǎn shí fēn} — Ten past eight.
\end{quote}

Its characters come later; for now it is a word for the ear.

\begin{cousinweb}
\emph{fēn} is inherited from Middle Chinese \emph{pjun} (Guangyun, in Baxter's transcription), glossed 'to divide'. Its use for minutes is not given there.
\end{cousinweb}

\begin{grammarlens}[title={What you've built}]
One more word for times and partings.
\end{grammarlens}

\subsection*{Your turn}
\begin{itemize}
\raggedright
  \item \emph{Say it:} \emph{fēn}
  \item \emph{Say it:} \emph{fēn}, once more
  \item \emph{Say it:} \emph{bā diǎn shí fēn} — Ten past eight.
  \item \emph{Recall:} say the Mandarin for date, then the Mandarin for o'clock, then say \emph{fēn} again
\end{itemize}

\begin{tcolorbox}[breakable,colback=teal!4,colframe=teal!35!black,title={Before you move on}]
What does fēn mean? (\textquotedblleft{}Minute; to divide\textquotedblright{}.) Say it once more.
\end{tcolorbox}
\section[bàn]{bàn — half}

\label{lesson:ZH-C48-ban}



\begin{quote}
Before the new one: say the Mandarin for o'clock, then the Mandarin for minute.
\end{quote}

\subsection*{You'll want to know: bàn}
\emph{bàn} — \textquotedblleft{}half\textquotedblright{}.

\begin{quote}
\emph{bā diǎn bàn} — Half past eight.
\end{quote}

Its characters come later; for now it is a word for the ear.

\begin{cousinweb}
\emph{bàn} is inherited from Middle Chinese \emph{panH} (Guangyun, in Baxter's transcription), glossed 'a thing divided in the middle'.
\end{cousinweb}

\begin{grammarlens}[title={What you've built}]
One more word for times and partings.
\end{grammarlens}

\subsection*{Your turn}
\begin{itemize}
\raggedright
  \item \emph{Say it:} \emph{bàn}
  \item \emph{Say it:} \emph{bàn}, once more
  \item \emph{Say it:} \emph{bā diǎn bàn} — Half past eight.
  \item \emph{Recall:} say the Mandarin for o'clock, then the Mandarin for minute, then say \emph{bàn} again
\end{itemize}

\begin{tcolorbox}[breakable,colback=teal!4,colframe=teal!35!black,title={Before you move on}]
What does bàn mean? (\textquotedblleft{}Half\textquotedblright{}.) Say it once more.
\end{tcolorbox}
\section[jǐ]{jǐ — how many; a few}

\label{lesson:ZH-C48-ji}



\begin{quote}
Before the new one: say the Mandarin for minute, then the Mandarin for half.
\end{quote}

\subsection*{You'll want to know: jǐ}
\emph{jǐ} — \textquotedblleft{}how many; a few\textquotedblright{}.

\begin{quote}
\emph{xiànzài jǐ diǎn} — What time is it?
\end{quote}

Its characters come later; for now it is a word for the ear.

\emph{jǐ} asks for a small number, as in \emph{jǐ diǎn}; \emph{duōshao} asks for any amount.

\begin{cousinweb}
\emph{jǐ} is inherited from Middle Chinese \emph{kj+jX} (Guangyun, in Baxter's transcription), glossed 'how many'.
\end{cousinweb}

\begin{grammarlens}[title={What you've built}]
One more word for times and partings.
\end{grammarlens}

\subsection*{Your turn}
\begin{itemize}
\raggedright
  \item \emph{Say it:} \emph{jǐ}
  \item \emph{Say it:} \emph{jǐ}, once more
  \item \emph{Say it:} \emph{xiànzài jǐ diǎn} — What time is it?
  \item \emph{Recall:} say the Mandarin for minute, then the Mandarin for half, then say \emph{jǐ} again
\end{itemize}

\begin{tcolorbox}[breakable,colback=teal!4,colframe=teal!35!black,title={Before you move on}]
What does jǐ mean? (\textquotedblleft{}How many; a few\textquotedblright{}.) Say it once more.
\end{tcolorbox}
\section[děng]{děng — to wait; to wait for}

\label{lesson:ZH-C48-deng}



\begin{quote}
Before the new one: say the Mandarin for half, then the Mandarin for how many.
\end{quote}

\subsection*{You'll want to know: děng}
\emph{děng} — \textquotedblleft{}to wait; to wait for\textquotedblright{}.

\begin{quote}
\emph{wǒ děng nǐ} — I'll wait for you.
\end{quote}

Its characters come later; for now it is a word for the ear.

\begin{cousinweb}
\emph{děng} is inherited from Middle Chinese \emph{tongX} (Guangyun, in Baxter's transcription), glossed 'equal, of a kind'. The sense 'to wait' is not given there, so how it got it is not settled.
\end{cousinweb}

\begin{grammarlens}[title={What you've built}]
One more word for times and partings.
\end{grammarlens}

\subsection*{Your turn}
\begin{itemize}
\raggedright
  \item \emph{Say it:} \emph{děng}
  \item \emph{Say it:} \emph{děng}, once more
  \item \emph{Say it:} \emph{wǒ děng nǐ} — I'll wait for you.
  \item \emph{Recall:} say the Mandarin for half, then the Mandarin for how many, then say \emph{děng} again
\end{itemize}

\begin{tcolorbox}[breakable,colback=teal!4,colframe=teal!35!black,title={Before you move on}]
What does děng mean? (\textquotedblleft{}To wait; to wait for\textquotedblright{}.) Say it once more.
\end{tcolorbox}
